% TOP_PROBLEM: {"id": 4600046, "title": "Periodic and finite configurations", "category_id": 11, "category": {"id": 11, "name": "dynamical_systems", "display_name": "Dynamical Systems", "description": "Problems about long-term behavior of deterministic systems, Hamiltonian dynamics, and periodic orbits.", "slug": "dynamical-systems", "order_index": 11, "created_at": "2026-07-31T15:26:25.671Z"}, "status": "partially_solved", "problem_number": "AMR-045-0046", "set_id": 13}
% FIRST_POSTED: 2026-10-01
% ENTRY_KIND: statement_only
% UnsolvedMath ID 4600046; source catalogue code AMR-045-0046
% !TeX program = lualatex
% Definition and sources only; this file is not an LLM solution attempt.
\documentclass[11pt,a4paper]{article}
\usepackage[margin=25mm]{geometry}
\usepackage{fontspec}
\setmainfont{lmroman10-regular.otf}[BoldFont=lmroman10-bold.otf,
 ItalicFont=lmroman10-italic.otf,BoldItalicFont=lmroman10-bolditalic.otf]
\usepackage{amsmath,amssymb,mathtools}
\usepackage{unicode-math}
\setmathfont{latinmodern-math.otf}
\AtBeginDocument{\let\setminus\smallsetminus}
\usepackage{xurl,hyperref}
\hypersetup{colorlinks=true,urlcolor=blue}
\setcounter{secnumdepth}{0}
\setlength{\parindent}{0pt}
\setlength{\parskip}{5pt}
\setlength{\emergencystretch}{3em}
\begin{document}
\section*{Periodic and finite configurations}\label{4600046}
{\small UnsolvedMath ID 4600046 · AMR-045-0046 · Dynamical Systems · AMR Open Problem Lists}
\subsection{Definitions and mathematical statement}
Let $d\ge2$, let $A$ be a finite alphabet with a distinguished
quiescent state $q$, and let $F:A^{\mathbb Z^d}\to A^{\mathbb Z^d}$
be a cellular automaton. Thus $F$ is continuous and commutes
with all translations; quiescence means $F(q^{\mathbb Z^d})=
q^{\mathbb Z^d}$. Define
\[
 P=\{x:\text{the translation orbit of }x\text{ is finite}\},
 \qquad C=\{x:\#\{v:x_v\ne q\}<\infty\}.
\]
Translation equivariance gives $F(P)\subset P$, and locality
and quiescence give $F(C)\subset C$. Write $F_P:P\to P$
and $F_C:C\to C$ for these restrictions.

Decide each of the following implications for arbitrary such
automata:
\[
 F_P\text{ injective}\ \Longrightarrow\ F_C\text{ surjective};
 \qquad
 F_C\text{ surjective}\ \Longrightarrow\ F_P\text{ surjective};
\]
\[
 F\text{ surjective}\ \Longrightarrow\ F_P\text{ surjective}.
\]
The last implication does not need a quiescent state, since it
does not involve finite-support configurations.

Spatial periodicity here means periodic in all lattice directions,
or equivalently fixed by a finite-index subgroup. Surjectivity
of $F_P$ permits a periodic preimage with a larger spatial
period than that of the target. Likewise, a finite-support
preimage may occupy a larger region than its target. Injectivity
and surjectivity refer to these entire sets, not to a single
finite box with imposed boundary conditions. Each implication
is a separate yes-or-no target.
\subsection{Short English statement}
Determine how injectivity and surjectivity on periodic configurations, finite-support configurations, and the full lattice constrain one another.
\subsection{Sources}
\begin{itemize}
\item[S1] ULAM.AI and the UnsolvedMath Contributors, supplied problems.json, record 4600046 (AMR-045-0046), AMR Open Problem Lists. Catalogue identity and source transcription; CC BY 4.0.
\url{https://huggingface.co/datasets/ulamai/UnsolvedMath}.
\item[S2] Boyle - Open Problems in Symbolic Dynamics (2008), Problem/Question/Conjecture 25.6. Original source indexed by AMR-045-0046.
\url{https://www.math.umd.edu/~mboyle/papers/openfinalsub3nov2007.pdf}.
\end{itemize}
\end{document}
